\section{Discussion and Limitations}
\label{sec:discussion}

The Qwen Effusion result changes the attribution question after steering succeeds. Its locked-dose response reproduces beyond random and sham controls, yet a fixed clinical alternative produces a larger change on the same patients and endpoint. Non-semantic controls establish an effective write direction; matched clinical competition tests whether its training label grants it relative privilege. The negative ownership interval answers the latter comparison without identifying which clinical concept the model uses.

Mass shows why this comparison also needs a specified prompt. Its advantage over all five clinical alternatives reproduces under both \emph{show} A/B instructions, while both \emph{is} margins are negative. Explicit \emph{show} yes/no has a similar observed margin, and the paired wording contrasts support a contribution from wording rather than answer letters alone. Clinical advantage and specificity remain distinct: larger effects among 119 random directions defeat the full specificity criteria.

These findings concern control of an answer, not improved diagnosis. Across Mass confirmation prompts, all AUROC-change intervals span zero and Brier error increases. The label-conditioned shifts in the separate ownership matrix cohort likewise give no basis for treating the average affirmative response as preferential correction in disease-positive patients (Tables~\ref{tab:label-conditioned-shifts} and~\ref{tab:label-conditioned-shifts-continued}). Under \emph{is}, clinical margins improve from the original anchor to explicit yes/no to A/B instructions, but that change in relative competition does not translate into demonstrated diagnostic benefit.

\paragraph{Limitations.}
The independent findings condition on ChestX-ray14, one Qwen model and visual locus, a fixed fitted direction bundle, and the evaluated doses and prompts. Relative clinical advantage depends on the quality and geometry of these fitted directions, while patient-bootstrap intervals quantify uncertainty over patients, not retraining of directions. The Mass wording comparison covers \emph{show} and \emph{is} within the evaluated instructions; it does not isolate a word-level mechanism from tokenization or prompt-prior interactions. Mass confirmation has 12 positives among 200 patients. Paired prompt contrasts are descriptive, and intervals spanning zero do not establish equivalence. The discovery crossover reused 200 patients and 20 random controls; its energy endpoint covers four calibration-eligible questions.

Report-derived labels, acquisition correlations, and mean pooling limit clinical interpretation. Selectivity against recurring-type random-label controls measures an advantage over those nuisance-label tasks; it does not rule out confounding by disease co-occurrence or acquisition features. Correlated findings, mixed representations, and shared answer-sensitive channels remain possible explanations for clinical cross-effects. The label-conditioned analysis describes differences in mean logit-margin shifts, leaving the disease conditionality of the underlying intervention mechanism unresolved.

The matrix, closure, and LLaVA analyses delimit the causal interpretation (Section~\ref{sec:mechanism-results}). They establish neither a shared answer mechanism nor reproducible reader--answer separation. Pilot intervals are unadjusted for selection or multiplicity, and the independent reader comparison uses a different sampling frame (Appendices~\ref{app:validation-readout} and~\ref{app:image-readout}). Metadata-matched longitudinal pairs retain unrecorded differences; their opportunity and direct model contrast remain unresolved. The cross-model comparison covers native configurations, including training and preprocessing, with half-credit ties under evaluation configured for BF16 (Appendix~\ref{app:paired-opportunity}).
